\section{Introduction}
\label{sec:intro}

High-quality paired training data is fundamental to supervised image editing models.
For 3D-aware editing tasks---such as rotating an object to a specified viewpoint,
changing its lighting, or altering its material---each training sample requires
a \emph{source image}, a \emph{target image} from a different viewpoint or condition,
and a corresponding \emph{editing instruction}.
Collecting such data from real-world captures is prohibitively expensive and laborious,
while existing synthetic approaches face a critical quality bottleneck.

\paragraph{The quality gap in synthetic data pipelines.}
A standard approach is to render 3D assets automatically, producing paired images at low cost.
However, na\"{i}ve rendering often yields scenes with artifacts:
flat lighting that fails to match the scene environment,
color inconsistencies between object and background,
physically implausible object placement (floating or intersecting the ground),
and poor shadow quality.
Existing pipelines typically address quality control through rigid scoring rules or
fixed controllers that select rendering adjustments based on numeric thresholds---approaches
that lack the semantic understanding needed to diagnose and resolve complex visual failures.

\paragraph{Our approach.}
We propose \aris{}, a pipeline in which rendering quality is improved through
a \emph{free-form VLM-guided evolution loop}.
Rather than relying on structured scores or fixed rules,
we employ a vision-language model (VLM) to review rendered scenes using natural language,
identifying specific visual problems such as flat lighting, color shift, or weak grounding.
An AI agent directly reads this free-form feedback and selects actions from a discrete
rendering action space to address the diagnosed issues---adjusting key light intensity,
environment rotation, material parameters, or object placement.
This process repeats until the VLM reviewer explicitly confirms the result is acceptable.

The AI agent replaces the manual trial-and-error that human artists would otherwise perform
when integrating 3D assets into Blender scenes:
once the pipeline is validated, large-scale dataset construction no longer requires
per-object expert tuning, reducing cost and improving consistency.

\paragraph{Instantiation: view-controlled rotation editing.}
We instantiate \aris{} for a concrete and evaluable task:
constructing a dataset for \emph{rotation-conditioned image editing},
where a model must rotate an object from a canonical front view to a target viewpoint
specified in natural language (e.g., \textit{``right side view''}, \textit{``back-left view''}).
We construct \textbf{ARIS-Rotate}, comprising 50 objects $\times$ 8 horizontal viewpoints,
yielding 350 editing pairs for training.

\paragraph{Results.}
Ablation experiments show that the VLM evolution loop substantially
improves rendering quality compared to single-pass rendering without iterative refinement.
LoRA fine-tuning of Qwen Image Edit 2511 on ARIS-Rotate improves over the base model
on PSNR, SSIM, and LPIPS, confirming that automatically constructed data
has downstream training value.
Comparing LoRA trained on refined data versus unrefined data further demonstrates
that the evolution loop produces data of higher utility.

\paragraph{Contributions.}
\begin{enumerate}
  \item \textbf{Free-form VLM-guided rendering evolution loop.}
    We introduce a quality refinement loop in which an AI agent reads
    free-form VLM reviewer feedback and selects rendering actions accordingly,
    replacing rigid score-based controllers with semantically-aware adaptation.

  \item \textbf{Scene-aware Blender rendering integration.}
    We describe a rendering system that preserves real-scene lighting and environment,
    integrating 3D objects into existing Blender scene files while allowing
    per-object adjustment of lighting, placement, material, and scene parameters
    through 24 structured atomic actions.

  \item \textbf{ARIS-Rotate dataset and validation.}
    We construct a 50-object, 350-pair rotation editing dataset with
    natural-language viewpoint prompts, and validate its training utility
    via LoRA fine-tuning on Qwen Image Edit 2511.
\end{enumerate}
